\ChapterOpener{29}{The Theory of Strategic Adjacencies}{Expansion Vectors, Diversification, and Option Value}{When and how should a successful core business expand into adjacent products, markets, or vertical integration?}

\begin{quote}
\textit{\textbf{Central thesis.} The best adjacency is not the nearest-looking market. It is an opportunity whose economics and execution become materially better because of capabilities the organization already possesses, and whose successful pursuit leaves those capabilities stronger rather than depleted. Expansion can create economies of scope and strategic options; it can also convert 1 focused, healthy business into 3 fragile, unrelated businesses.}
\end{quote}

\noindent\textbf{Reader's note and orientation.} This chapter examines how enterprises expand their boundaries without losing the capabilities that make them valuable. It provides a founder's guide to product extension, market entry, vertical integration, diversification, portfolio strategy, acquisition, real options, and the disciplined refusal of attractive distractions. It builds on Chapter 27 (scaling) and Chapter 28 (competitive moats), establishing the foundation for Chapter 30 (ecosystems and strategic partnerships) and Chapter 31 (international expansion).

\vspace{0.5em}
\noindent\textbf{Scope and evidence convention.} Published theory and cited research are identified in Appendix D. Frameworks labeled \textit{proposed} are synthesis, not validated scientific scales. All commercial figures in this chapter are hypothetical exercises, not actual Bayes Institute revenue, contracts, or market data. Legal and accounting material is educational, not jurisdiction-specific professional advice.

\begin{table*}[t]
\centering
\small
\caption{Reading Map: The Theory of Strategic Adjacencies}
\label{tab:ch29_reading_map}
\begin{tabularx}{\textwidth}{l X}
\toprule
\textbf{Part} & \textbf{Primary Question} \\
\midrule
I. The strategic problem of going somewhere else & What makes expansion to a 2nd opportunity more dangerous than starting the 1st? \\
II. Mapping the adjacency space & How do we measure the multi-dimensional distance between capabilities and markets? \\
III. The economics of scope, synergy, and attention & When do shared assets create genuine scope economies rather than congestion taxes? \\
IV. Vertical integration and the boundaries of the firm & When should an enterprise make, buy, or partner along its critical value chain? \\
V. Diversification and portfolio strategy & How can a multi-business portfolio balance cash cows, growth engines, and options? \\
VI. Acquisitions, joint ventures, and corporate development & When does buying an asset destroy more value than organic staging or alliances? \\
VII. The industry atlas: where expansion mechanisms differ & How do adjacency mechanics vary across 12 distinct industry architectures? \\
VIII. Bayes Institute: an evolutionary adjacency portfolio & How should an early STEM institute sequence cohorts, colleges, enterprise, and research? \\
IX. The founder's adjacency decision system & What concrete screening memo, gating rules, and worksheets govern expansion choices? \\
X. Advanced decision laboratories & What analytical traps make unprofitable extensions look profitable on paper? \\
Practical appendices and research foundations & What checklists, break-even tools, and foundational scholarship guide execution? \\
\bottomrule
\end{tabularx}
\end{table*}

\subimport{./}{00-the-strategic-problem-of-going-somewhere-else.tex}
\subimport{./}{01-mapping-the-adjacency-space.tex}
\subimport{./}{02-economics-of-scope-synergy-and-attention.tex}
\subimport{./}{03-vertical-integration-and-boundaries-of-the-firm.tex}
\subimport{./}{04-diversification-and-portfolio-strategy.tex}
\subimport{./}{05-acquisitions-joint-ventures-corporate-development.tex}
\subimport{./}{06-the-industry-atlas.tex}
\subimport{./}{07-bayes-institute-evolutionary-portfolio.tex}
\subimport{./}{08-founders-adjacency-decision-system.tex}
\subimport{./}{09-advanced-decision-laboratories.tex}
\subimport{./}{10-practical-appendices-and-foundations.tex}
